\chapter{Conclusion}\label{ch:conclusion}

This dissertation moved the biomimetic humanoid robot project from an unvalidated design concept to a quantitatively characterized actuator set and a validated knee-level biomechanical design, and it defined the path from those experimental foundations to neuromechanical control.

At the actuator level, the analysis in Chapters~\ref{ch:methods} and~\ref{ch:results} produced novel equations for calculating isometric force in Festo \qtylist{10;20}{\mm} diameter BPAs. The work elucidated the relationship between maximum BPA force at \qty{620}{\kPa} and resting length, $F_{620}(l_0)$, a previously unreported characteristic of these actuators that the manufacturer's tool predicts with the wrong trend, and it produced a nondimensional force surface, $F^*(\epsilon^*,P^*)$, with only three diameter-dependent coefficients. Taken together, these fits predict isometric BPA force across resting lengths from a small number of measurements, correct the large over-predictions of the manufacturer's tool at short resting lengths, and outperform published single-length models. In a design workflow, these equations make it possible to select actuator diameter and resting length, and to evaluate whether a proposed muscle routing can meet force and range-of-motion requirements, before committing to a prototype.

At the joint level, the improved actuator model was combined with a compliance- and wrapping-corrected musculotendon kinematics model and validated on two knee designs. The hybrid torque calculation method isolated the sources of error between prediction and measurement, and the resulting correction terms --- a constant length offset, series compliance represented by an equivalent stiffness, and a wrapping loss term --- substantially improved torque prediction on the pinned knee and transferred to the biomimetic four-bar knee. The biomimetic knee driven by optimized BPAs was shown to meet or exceed human isometric knee torque over most of the range of motion, confirming that the optimized muscle routing of the author's earlier work is realizable in hardware, and identifying the high-flexion regime where additional flexor capacity is needed.

Beyond the specific models, the dissertation's methodological contribution is the actuator-to-joint design chain itself: characterize the actuator, correct for the compliance and geometry of the transmission, validate against human biomechanics, and only then close the loop with a controller. The remaining gap --- the neural layer --- is scoped in Chapter~\ref{ch:futurework}: a MuJoCo plant converted from the lab's OpenSim models via MyoConverter, wrapped in an SNS-Toolbox synthetic nervous system implementing a two-layer CPG with per-leg rhythm generation and pattern formation, left--right coordination, and Ia, Ib, II, and contact feedback, verified first by tonic-stimulation tests against the literature, then extended with vestibular, ocular, and cerebellar layers, and finally deployed on embedded hardware controlling a treadmill-walking robot. The biomimetic platform that results will permit neuromechanical experiments --- neural lesions, sensory deletions, pathway stimulation --- that are impossible in living animals, closing the cycle between better robots and better biological understanding, and opening the way toward dynamic whole-body maneuvers.
